%-----------------------------------
% KI-Verzeichnis
% Leitfaden Abschnitt 6.2.4: "Im KI-Verzeichnis, das dem Literaturverzeichnis
% folgt, sind die Referenzen entsprechend aufzulisten." Format:
%   Open AI. (2023). ChatGPT (Version vom 14. Maerz) [Grosses Sprachmodell].
%   (https://chat.openai.com/chat)
%
% ZU PRUEFEN: Die Liste enthaelt nur das Werkzeug, dessen Einsatz beim
% Aufbereiten dieser Datei bekannt ist. Jedes weitere fuer Textpassagen,
% Abbildungen oder Recherche genutzte Tool ist zu ergaenzen. Wurde keine KI
% verwendet, sind dieses Verzeichnis und Anhang II zu entfernen.
%-----------------------------------
\newpage
\phantomsection
\addcontentsline{toc}{section}{\langde{KI-Verzeichnis}\langen{List of AI tools}}
\section*{\langde{KI-Verzeichnis}\langen{List of AI tools}}

\begingroup
\setlength{\parindent}{0pt}
\setlength{\parskip}{0.7em}

\hangindent=1.5em\hangafter=1
Anthropic. (2026). Claude (Version Opus~5, Abfrage vom 30.~Juli 2026) [Gro{\ss}es Sprachmodell]. (\url{https://claude.ai})

\hangindent=1.5em\hangafter=1
$\bullet$ (weitere verwendete Werkzeuge hier erg\"anzen: Herausgeber, Jahr, Bezeichnung, Version, Datum der Abfrage, URL)

\endgroup
